\begin{titlepage}

\projectAuthorNameLT. \projectTitleLT. \projectTypeLT / supervisor \projectSupervisorNameFirstLT{}; advisor \projectSupervisorNameSecondLT{}; Kaunas University of Technology, \projectFacultyLT.

Study field and area (study field group): \projectStudyFieldAndAreaLT.

Keywords: \projectKeywordsLT.

\projectCityLT, \projectYearLT. \pageref{LastPage}.


\begin{center}
  \textbf{Santrauka}
\end{center}

This literature analysis examines how functional near-infrared spectroscopy (fNIRS) and machine learning can support objective assessment of neurocognitive and mental-health related indicators. It does not treat the current evidence as sufficient for stand-alone diagnosis. The thesis position is limited to relative haemodynamic task patterns, not absolute tissue oxygen saturation or disorder-specific diagnosis.

The revised report focuses on the methodological safeguards needed for defensible fNIRS machine learning. It emphasizes short-channel regression for extracerebral systemic activity, preservation of channel-by-time signal structure, correct interpretation of \(\Delta\mathrm{HbO}\), \(\Delta\mathrm{HbR}\), HbDiff, and HbT, and leakage-safe participant-wise evaluation. These safeguards reduce the risk that a model learns scalp blood flow, heart rate, participant identity, or preprocessing artefacts instead of task-related haemodynamic patterns.




\end{titlepage}
